\documentclass[aspectratio=169,11pt]{beamer}
\usetheme{Madrid}
\usecolortheme{dolphin}
\setbeamertemplate{navigation symbols}{}
\setbeamertemplate{caption}[numbered]

\usepackage[utf8]{inputenc}
\usepackage[T1]{fontenc}
\usepackage[spanish,es-noshorthands]{babel}
\usepackage{amsmath,amssymb,amsthm,mathtools}
\usepackage{tikz}
\usepackage{pgfplots}
\pgfplotsset{compat=1.18}
\usepackage{booktabs}
\usepackage{array}

\title[Prob. y Estadística]{Clase 28: Bondad de ajuste, independencia y homogeneidad}
\subtitle{Unidad 7: Prueba de hipótesis}
\institute[UNS]{Universidad Nacional del Sur \\ Departamento de Matemática}
\author{Probabilidad y Estadística (5765)}
\date{}

\begin{document}

\begin{frame}
\titlepage
\end{frame}

\begin{frame}{Contenidos}
\tableofcontents
\end{frame}

\section{Introducción}

\begin{frame}{Pruebas chi-cuadrado}
Las pruebas vistas hasta ahora se aplican a parámetros de variables cuantitativas (medias, varianzas). Ahora estudiamos una familia de pruebas basadas en la comparación de \textbf{frecuencias observadas} contra \textbf{frecuencias esperadas}, aplicables a datos categóricos (de conteo).

\begin{block}{Estadístico general de Pearson}
En las tres pruebas de esta clase se utiliza el mismo tipo de estadístico:
\[
\chi^2 = \sum_{i} \frac{(O_i-E_i)^2}{E_i},
\]
donde $O_i$ son las frecuencias \textbf{observadas} y $E_i$ las frecuencias \textbf{esperadas} bajo $H_0$, en cada una de las categorías o celdas consideradas.
\end{block}

\begin{itemize}
\item \textbf{Bondad de ajuste}: ¿los datos provienen de una distribución teórica dada?
\item \textbf{Independencia}: ¿dos variables categóricas, medidas sobre la misma población, están relacionadas?
\item \textbf{Homogeneidad}: ¿varias poblaciones tienen la misma distribución para una variable categórica?
\end{itemize}
\end{frame}

\begin{frame}{Región crítica de las pruebas $\chi^2$}
\centering
\begin{tikzpicture}
\begin{axis}[
  width=9.5cm, height=5cm,
  axis lines=left, ytick=\empty,
  xlabel={$\chi^2$}, ylabel={densidad},
  xmin=0, xmax=25, ymin=0,
]
\addplot[domain=0.1:25, samples=150, thick, blue, name path=curve] {x^(1.5)*exp(-x/2)};
\addplot[domain=11.07:25, samples=60, fill=red!40, draw=none] {x^(1.5)*exp(-x/2)} \closedcycle;
\draw[dashed] (axis cs:11.07,0) -- (axis cs:11.07,0.85);
\node at (axis cs:11.07,-0.15) {\footnotesize $\chi^2_{k-1,\alpha}$};
\end{axis}
\end{tikzpicture}
\vspace{0.2cm}

En las tres pruebas de esta clase, la región crítica está \textbf{siempre en la cola derecha}: valores grandes de $\chi^2$ (mucha discrepancia entre lo observado y lo esperado) llevan a rechazar $H_0$.
\end{frame}

\section{Bondad de ajuste}

\begin{frame}{Prueba de bondad de ajuste}
Se dispone de $n$ observaciones clasificadas en $k$ categorías mutuamente excluyentes, con frecuencias observadas $O_1,\dots,O_k$ ($\sum O_i=n$). Se desea probar si los datos son compatibles con una distribución teórica que asigna probabilidades $p_1,\dots,p_k$ a cada categoría:
\[
H_0: P(\text{categoría } i) = p_i, \ i=1,\dots,k \qquad \text{vs.} \qquad H_1: \text{para algún } i,\ P(\text{categoría } i)\neq p_i.
\]

\begin{block}{Frecuencias esperadas y estadístico}
\[
E_i = n\,p_i, \qquad \chi^2 = \sum_{i=1}^{k}\frac{(O_i-E_i)^2}{E_i} \ \overset{H_0}{\approx}\ \chi^2_{k-1-m},
\]
donde $m$ es la cantidad de parámetros de la distribución teórica que debieron \textbf{estimarse} a partir de los datos (si la distribución está completamente especificada, $m=0$).
\end{block}

\textbf{Región crítica:} $\chi^2 > \chi^2_{k-1-m,\,\alpha}$ (siempre cola derecha: valores grandes de $\chi^2$ indican mal ajuste).
\end{frame}

\begin{frame}{Condiciones de aplicación}
\begin{alertblock}{Requisito sobre las frecuencias esperadas}
La aproximación de la distribución de $\chi^2$ es adecuada cuando las frecuencias esperadas no son demasiado pequeñas. Criterio usual:
\[
E_i \ge 5 \quad \text{para todas las categorías (o al menos el $80\%$ de ellas).}
\]
Si alguna $E_i<5$, se recomienda \textbf{agrupar categorías} contiguas hasta cumplir la condición.
\end{alertblock}
\end{frame}

\begin{frame}{Ejemplo: ajuste a una distribución uniforme}
\begin{example}
Se lanza un dado $120$ veces para verificar si está equilibrado. Frecuencias observadas:
\begin{center}
\begin{tabular}{@{}lcccccc@{}}
\toprule
Cara & 1 & 2 & 3 & 4 & 5 & 6 \\
\midrule
$O_i$ & 25 & 17 & 15 & 23 & 24 & 16 \\
\bottomrule
\end{tabular}
\end{center}
Con $\alpha=0.05$, pruebe si el dado es equilibrado.

\textbf{1. Hipótesis:} $H_0: p_i=1/6\ \forall i$ vs. $H_1:$ algún $p_i\neq1/6$.

\textbf{2. Frecuencias esperadas:} $E_i=120\times\tfrac{1}{6}=20$ para todo $i$ (todas $\ge5$, se cumple la condición).

\textbf{3. Estadístico y gl:} $k=6$, $m=0 \Rightarrow$ gl$=5$.

\textbf{4. Región crítica:} $\chi^2>\chi^2_{5,\,0.05}=11.07$.

\textbf{5. Cálculo:}
\[
\chi^2=\frac{5^2}{20}+\frac{3^2}{20}+\frac{5^2}{20}+\frac{3^2}{20}+\frac{4^2}{20}+\frac{4^2}{20}=\frac{25+9+25+9+16+16}{20}=\frac{100}{20}=5.0.
\]

\textbf{6. Decisión y conclusión:} $5.0<11.07\Rightarrow$ no se rechaza $H_0$; no hay evidencia de que el dado esté desequilibrado.
\end{example}
\end{frame}

\begin{frame}{Ejemplo: ajuste a una distribución con parámetro estimado}
\begin{example}
Se registra el número de defectos por unidad en $200$ piezas, y se desea comprobar si sigue una distribución de Poisson (la media $\lambda$ debe estimarse de los datos: $\hat\lambda=1.5$).
\begin{center}
\small
\begin{tabular}{@{}lccccc@{}}
\toprule
N° defectos & 0 & 1 & 2 & 3 & $\ge4$ \\
\midrule
$O_i$ & 48 & 65 & 48 & 26 & 13 \\
$p_i$ (Poisson, $\hat\lambda=1.5$) & 0.223 & 0.335 & 0.251 & 0.126 & 0.065 \\
$E_i=200\,p_i$ & 44.6 & 67.0 & 50.2 & 25.2 & 13.0 \\
\bottomrule
\end{tabular}
\end{center}
Con $\alpha=0.05$: gl $=k-1-m=5-1-1=3$ (se estimó $\lambda$, $m=1$).

\textbf{Región crítica:} $\chi^2>\chi^2_{3,\,0.05}=7.815$.

\textbf{Cálculo:}
\[
\chi^2=\frac{(48-44.6)^2}{44.6}+\frac{(65-67.0)^2}{67.0}+\frac{(48-50.2)^2}{50.2}+\frac{(26-25.2)^2}{25.2}+\frac{(13-13.0)^2}{13.0}\approx0.259+0.060+0.096+0.025+0.000=0.44.
\]

\textbf{Decisión:} $0.44<7.815\Rightarrow$ no se rechaza $H_0$: el ajuste a la Poisson es adecuado.
\end{example}
\end{frame}

\section{Prueba de independencia}

\begin{frame}{Prueba de independencia (tablas de contingencia)}
Se clasifica una muestra de $n$ individuos simultáneamente según \textbf{dos variables categóricas}, $A$ (con $r$ categorías) y $B$ (con $c$ categorías), formando una tabla de contingencia $r\times c$ de frecuencias observadas $O_{ij}$.

\begin{block}{Hipótesis}
\[
H_0: \text{las variables } A \text{ y } B \text{ son independientes} \qquad \text{vs.} \qquad H_1: \text{no son independientes}.
\]
\end{block}

\begin{block}{Frecuencias esperadas bajo $H_0$}
Si $O_{i\bullet}$ es el total de la fila $i$, $O_{\bullet j}$ el total de la columna $j$, y $n$ el total general:
\[
E_{ij} = \frac{O_{i\bullet}\cdot O_{\bullet j}}{n}.
\]
\end{block}

\begin{block}{Estadístico}
\[
\chi^2=\sum_{i=1}^{r}\sum_{j=1}^{c}\frac{(O_{ij}-E_{ij})^2}{E_{ij}} \ \overset{H_0}{\approx}\ \chi^2_{(r-1)(c-1)}.
\]
\end{block}
\textbf{Región crítica:} $\chi^2>\chi^2_{(r-1)(c-1),\,\alpha}$.
\end{frame}

\begin{frame}{Ejemplo: independencia entre turno y ausentismo}
\begin{example}
Se estudia si el ausentismo laboral es independiente del turno de trabajo, sobre $300$ empleados:
\begin{center}
\begin{tabular}{@{}lccc|c@{}}
\toprule
 & Mañana & Tarde & Noche & Total \\
\midrule
Con ausentismo & 18 & 22 & 35 & 75 \\
Sin ausentismo & 82 & 78 & 65 & 225 \\
\midrule
Total & 100 & 100 & 100 & 300 \\
\bottomrule
\end{tabular}
\end{center}
Con $\alpha=0.05$, ¿el ausentismo depende del turno?

\textbf{1. Hipótesis:} $H_0:$ ausentismo y turno son independientes vs. $H_1:$ no lo son.

\textbf{2. Frecuencias esperadas} ($E_{ij}=O_{i\bullet}O_{\bullet j}/300$):
\[
E_{11}=\frac{75\times100}{300}=25, \quad E_{12}=25, \quad E_{13}=25, \quad E_{21}=75, \quad E_{22}=75, \quad E_{23}=75.
\]
\end{example}
\end{frame}

\begin{frame}{Ejemplo: independencia (continuación)}
\begin{example}
\textbf{3. Grados de libertad:} $(r-1)(c-1)=(2-1)(3-1)=2$.

\textbf{4. Región crítica:} $\chi^2>\chi^2_{2,\,0.05}=5.991$.

\textbf{5. Cálculo:}
\[
\chi^2=\frac{(18-25)^2}{25}+\frac{(22-25)^2}{25}+\frac{(35-25)^2}{25}+\frac{(82-75)^2}{75}+\frac{(78-75)^2}{75}+\frac{(65-75)^2}{75}
\]
\[
=\frac{49}{25}+\frac{9}{25}+\frac{100}{25}+\frac{49}{75}+\frac{9}{75}+\frac{100}{75} = 1.96+0.36+4.00+0.653+0.12+1.333 = 8.43.
\]

\textbf{6. Decisión:} $8.43>5.991\Rightarrow$ se \textbf{rechaza} $H_0$.

\textbf{7. Conclusión:} al $5\%$ de significación, hay evidencia de que el ausentismo \textbf{no} es independiente del turno de trabajo (el turno noche muestra mayor ausentismo que el esperado bajo independencia).
\end{example}
\end{frame}

\section{Prueba de homogeneidad}

\begin{frame}{Prueba de homogeneidad}
Se toman muestras \textbf{independientes} de tamaño fijo de $r$ poblaciones distintas, y se clasifica cada observación según una variable categórica con $c$ categorías. Se desea probar si las $r$ poblaciones tienen la \textbf{misma distribución} para dicha variable.

\begin{block}{Hipótesis}
\[
H_0: \text{las } r \text{ poblaciones tienen la misma distribución en la variable categórica.}
\]
\[
H_1: \text{al menos una población difiere de las demás.}
\]
\end{block}

\begin{alertblock}{Diferencia clave con la prueba de independencia}
En \textbf{independencia}, se toma \textbf{una única muestra} de $n$ individuos y se los clasifica según dos variables (los totales marginales de fila y columna son ambos aleatorios). En \textbf{homogeneidad}, los tamaños de cada muestra (una por población, es decir, los totales de fila o columna) están \textbf{fijados de antemano} por el diseño. Sin embargo, el cálculo —frecuencias esperadas, estadístico $\chi^2$ y grados de libertad $(r-1)(c-1)$— es \textbf{formalmente idéntico}.
\end{alertblock}
\end{frame}

\begin{frame}{Ejemplo: homogeneidad entre proveedores}
\begin{example}
Se extraen muestras de $150$ piezas de cada uno de tres proveedores y se clasifican según su calidad:
\begin{center}
\begin{tabular}{@{}lccc|c@{}}
\toprule
 & Proveedor A & Proveedor B & Proveedor C & Total \\
\midrule
Primera calidad & 120 & 108 & 96 & 324 \\
Segunda calidad & 30 & 42 & 54 & 126 \\
\midrule
Total & 150 & 150 & 150 & 450 \\
\bottomrule
\end{tabular}
\end{center}
Con $\alpha=0.01$, ¿los tres proveedores son homogéneos respecto de la calidad entregada?

\textbf{1. Hipótesis:} $H_0:$ los tres proveedores tienen igual distribución de calidad vs. $H_1:$ al menos uno difiere.

\textbf{2. Frecuencias esperadas:} $E_{1j}=\dfrac{324\times150}{450}=108$ para cada proveedor (primera calidad); $E_{2j}=\dfrac{126\times150}{450}=42$ (segunda calidad).
\end{example}
\end{frame}

\begin{frame}{Ejemplo: homogeneidad (continuación)}
\begin{example}
\textbf{3. Grados de libertad:} $(r-1)(c-1)=(2-1)(3-1)=2$.

\textbf{4. Región crítica:} $\chi^2>\chi^2_{2,\,0.01}=9.210$.

\textbf{5. Cálculo:}
\[
\chi^2=\frac{(120-108)^2}{108}+\frac{(108-108)^2}{108}+\frac{(96-108)^2}{108}+\frac{(30-42)^2}{42}+\frac{(42-42)^2}{42}+\frac{(54-42)^2}{42}
\]
\[
=\frac{144}{108}+0+\frac{144}{108}+\frac{144}{42}+0+\frac{144}{42}=1.333+1.333+3.429+3.429=9.52.
\]

\textbf{6. Decisión:} $9.52>9.210\Rightarrow$ se \textbf{rechaza} $H_0$.

\textbf{7. Conclusión:} al $1\%$ de significación, hay evidencia de que los tres proveedores no son homogéneos en la calidad de las piezas entregadas (el proveedor C presenta peor calidad relativa).
\end{example}
\end{frame}

\begin{frame}{Bondad de ajuste, independencia y homogeneidad: comparación}
\begin{table}[]
\centering
\scriptsize
\begin{tabular}{@{}p{2.3cm}p{3.4cm}p{3.5cm}@{}}
\toprule
Prueba & Diseño de muestreo & Pregunta que responde \\
\midrule
Bondad de ajuste & una muestra, una variable categórica & ¿la variable sigue una distribución teórica dada? \\[0.6em]
Independencia & una muestra, dos variables categóricas & ¿las dos variables están relacionadas? \\[0.6em]
Homogeneidad & varias muestras (una por población), una variable categórica & ¿las poblaciones tienen la misma distribución? \\
\bottomrule
\end{tabular}
\end{table}
\begin{alertblock}{Punto en común}
Pese a responder preguntas distintas, las tres pruebas de independencia y homogeneidad comparten \textbf{exactamente} el mismo cálculo: $E_{ij}=O_{i\bullet}O_{\bullet j}/n$, $\chi^2=\sum(O_{ij}-E_{ij})^2/E_{ij}$, con $(r-1)(c-1)$ grados de libertad.
\end{alertblock}
\end{frame}

\section{Condiciones de aplicación generales}

\begin{frame}{Condiciones de aplicación de las pruebas $\chi^2$}
\begin{itemize}
\item Las observaciones deben ser \textbf{independientes} entre sí.
\item Cada individuo debe clasificarse en \textbf{una y solo una} categoría o celda.
\item Las frecuencias esperadas deben ser suficientemente grandes: criterio habitual $E_i\ge5$ (o $E_{ij}\ge5$) en todas las celdas, o al menos en el $80\%$ de ellas, sin celdas con $E_i<1$.
\item Si no se cumple, se deben \textbf{agrupar} categorías o filas/columnas adyacentes hasta satisfacer la condición, a costa de perder grados de libertad.
\item El estadístico $\chi^2$ de Pearson es siempre una prueba de \textbf{cola derecha}: valores grandes indican discrepancia entre lo observado y lo esperado bajo $H_0$.
\end{itemize}
\end{frame}

\section{Resumen}

\begin{frame}{Resumen de la clase}
\begin{table}[]
\centering
\scriptsize
\begin{tabular}{@{}p{2.6cm}p{4.6cm}p{2.3cm}@{}}
\toprule
Prueba & Hipótesis nula & Grados de libertad \\
\midrule
Bondad de ajuste & los datos siguen la distribución especificada & $k-1-m$ \\
Independencia & dos variables categóricas son independientes & $(r-1)(c-1)$ \\
Homogeneidad & varias poblaciones tienen igual distribución & $(r-1)(c-1)$ \\
\bottomrule
\end{tabular}
\end{table}

\begin{itemize}
\item Las tres pruebas comparten el estadístico $\displaystyle\chi^2=\sum\frac{(O-E)^2}{E}$ y la misma lógica de decisión: rechazar $H_0$ si $\chi^2$ es grande.
\item Se requieren frecuencias esperadas $E_i\ge5$ para que la aproximación $\chi^2$ sea válida; de lo contrario, agrupar categorías.
\item Con esta clase se completa la Unidad 7: quedan cubiertas las pruebas de hipótesis para parámetros de una y dos poblaciones, y las pruebas de ajuste, independencia y homogeneidad basadas en $\chi^2$.
\end{itemize}
\end{frame}

\end{document}
